\documentclass[12pt,a4paper]{article}

\usepackage[T1]{fontenc}
\usepackage[utf8]{inputenc}
\usepackage[romanian]{babel}
\usepackage{geometry}
\usepackage{listings}
\usepackage{xcolor}
\usepackage[hidelinks]{hyperref}

\geometry{margin=2.5cm}

\definecolor{codegreen}{rgb}{0,0.6,0}
\definecolor{codegray}{rgb}{0.5,0.5,0.5}

\lstset{
    language=C++,
    basicstyle=\ttfamily\small,
    keywordstyle=\color{blue},
    commentstyle=\color{codegreen},
    numberstyle=\tiny\color{codegray},
    numbers=left,
    frame=single,
    breaklines=true
}

\title{\textbf{Tutoriat 1}\\Încapsulare, Getters \& Setters, Constructori}
\author{}
\date{}

\begin{document}

\maketitle

\section{Programarea Orientată pe Obiecte}

Programarea Orientată pe Obiecte (POO) reprezintă o paradigmă de programare care organizează datele și funcționalitățile sub forma unor obiecte.

În loc să avem variabile și funcții separate:

\begin{lstlisting}
string nume;
int varsta;

void afiseazaNume() {
    cout << nume;
}
\end{lstlisting}

putem grupa aceste informații într-o singură structură numită clasă.

\section{Class vs Struct}

Clasele reprezintă o extensie a structurilor de date.

La fel ca structurile, ele pot conține variabile:

\begin{lstlisting}
class Student {
    int varsta;
    string nume;
};
\end{lstlisting}

Dar pot conține și funcții:

\begin{lstlisting}
class Student {
    int varsta;
    string nume;

    void afiseazaNume() {
        cout << nume;
    }
};
\end{lstlisting}

Datele și funcțiile sunt astfel grupate într-o singură entitate.

\section{Încapsularea}

Încapsularea este unul dintre principiile fundamentale ale programării orientate pe obiecte.

Ea se bazează pe două idei:

\begin{itemize}
    \item gruparea datelor și funcțiilor într-o singură structură;
    \item controlul accesului la date.
\end{itemize}

Pentru a controla accesul folosim modificatori de acces.

\section{Modificatori de acces}

Există trei modificatori de acces în C++.

\subsection*{private}

Membrii sunt accesibili doar din interiorul clasei.

\begin{lstlisting}
class Student {
private:
    int varsta;
};
\end{lstlisting}

\subsection*{public}

Membrii pot fi accesați de oriunde.

\begin{lstlisting}
class Student {
public:
    void afiseaza() {
        cout << "Salut!";
    }
};
\end{lstlisting}

\subsection*{protected}

Pentru moment îl putem considera asemănător cu \texttt{private}. Diferențele vor apărea atunci când vom discuta despre moștenire.


\newpage
\section{Sintaxa unei clase}

\begin{lstlisting}
class NumeClasa {

private:
    // date

public:
    // functii

};
\end{lstlisting}

Exemplu:

\begin{lstlisting}
class Student {

private:
    string nume;

public:
    void afiseazaNume() {
        cout << nume;
    }

};
\end{lstlisting}

Utilizatorul clasei poate apela:

\begin{lstlisting}
student.afiseazaNume();
\end{lstlisting}

dar nu poate accesa:

\begin{lstlisting}
student.nume;
\end{lstlisting}

\section{Getters și Setters}

Dacă datele sunt private, avem nevoie de metode prin care să le accesăm și să le modificăm.

\subsection*{Getter}

Un getter întoarce valoarea unei variabile private.

\begin{lstlisting}
int getVarsta() {
    return varsta;
}
\end{lstlisting}

\subsection*{Setter}

Un setter modifică valoarea unei variabile private.

\begin{lstlisting}
void setVarsta(int v) {
    varsta = v;
}
\end{lstlisting}

\subsection*{Exemplu}

\begin{lstlisting}
class Student {

private:
    string nume;
    int varsta;

public:

    void setNume(string n) {
        nume = n;
    }

    string getNume() {
        return nume;
    }

    void setVarsta(int v) {
        varsta = v;
    }

    int getVarsta() {
        return varsta;
    }
};
\end{lstlisting}

Utilizare:

\begin{lstlisting}
Student s;

s.setNume("Ion");
s.setVarsta(20);

cout << s.getNume();
cout << s.getVarsta();
\end{lstlisting}

\section{Sunt getter-ele întotdeauna utile?}

La prima vedere:

\begin{lstlisting}
student.varsta
\end{lstlisting}

și

\begin{lstlisting}
student.getVarsta()
\end{lstlisting}

par foarte asemănătoare.

Dacă un getter doar returnează valoarea unei variabile și nu face altceva, diferența practică este foarte mică.

\begin{lstlisting}
class Student {
private:
    int varsta;

public:
    int getVarsta() {
        return varsta;
    }
};
\end{lstlisting}

față de

\begin{lstlisting}
class Student {
public:
    int varsta;
};
\end{lstlisting}

obținem aproape același comportament, doar că scriem:

\begin{lstlisting}
student.getVarsta()
\end{lstlisting}

în loc de

\begin{lstlisting}
student.varsta
\end{lstlisting}

În practică este adesea mai util să expunem comportamente și acțiuni decât să expunem date.

\section{Obiecte}

\subsection*{Definiție formală}

Un obiect este o instanță a unei clase.

\subsection*{Definiție informală}

Un obiect este o variabilă al cărei tip este o clasă.

Exemplu:

\begin{lstlisting}
Student s;
\end{lstlisting}

Aici:

\begin{itemize}
    \item Student este clasa;
    \item s este obiectul.
\end{itemize}

\section{Constructori}

Până acum am creat obiecte și ulterior le-am setat valorile.

\begin{lstlisting}
Student s;

s.setNume("Ion");
s.setVarsta(20);
\end{lstlisting}

Putem simplifica procesul folosind un constructor.

\subsection*{Definiție}

Constructorul este o metodă specială:

\begin{itemize}
    \item nu are tip de retur;
    \item are același nume ca și clasa;
    \item este apelată automat la crearea obiectului.
\end{itemize}

\subsection*{Exemplu}

\begin{lstlisting}
class Student {

private:
    string nume;
    int varsta;

public:

    Student(string n, int v) {
        nume = n;
        varsta = v;
    }

    string getNume() {
        return nume;
    }

    int getVarsta() {
        return varsta;
    }
};
\end{lstlisting}

Utilizare:

\begin{lstlisting}
Student s("Ion", 20);

cout << s.getNume();
cout << s.getVarsta();
\end{lstlisting}

\newpage

\section{Exemplu complet}

Până acum am folosit getters și setters pentru a accesa date private.

Un student observă rapid că:

\begin{lstlisting}
card.getBani()
\end{lstlisting}

și

\begin{lstlisting}
card.bani
\end{lstlisting}

par să facă aproximativ același lucru.

În exemplul următor nu vom expune deloc suma de bani printr-un getter. În schimb vom lăsa obiectul să decidă singur dacă poate efectua o cumpărătură.

\subsection{Clasa Produs}

\begin{lstlisting}
class Produs {

private:
    string nume;
    int pret;

public:

    Produs(string n, int p)
        : nume(n), pret(p) {}

    string getNume() {
        return nume;
    }

    int getPret() {
        return pret;
    }
};
\end{lstlisting}

\newpage
\subsection{Clasa Card}

\begin{lstlisting}
class Card {

private:
    string posesor;
    int bani;

public:

    Card(string p, int b)
        : posesor(p), bani(b) {}

    bool buy(const Produs& produs) {

        if(bani >= produs.getPret()) {
            bani -= produs.getPret();
            return true;
        }

        return false;
    }

    void afiseazaSold() {

        cout << posesor
             << " are "
             << bani
             << " lei.\n";
    }
};
\end{lstlisting}

Observați că nu avem:

\begin{lstlisting}
int getBani();
\end{lstlisting}

Nu întrebăm cardul câți bani are pentru a decide noi dacă se poate cumpăra produsul.

În schimb scriem:

\begin{lstlisting}
card.buy(produs)
\end{lstlisting}

și obiectul decide singur dacă tranzacția poate avea loc.

\newpage
\subsection{Clasa Persoana}

\begin{lstlisting}
#include <vector>

class Persoana {

private:
    string nume;
    Card card;
    vector<Produs> produse;

public:

    Persoana(string n, Card c)
        : nume(n), card(c) {}

    void cumpara(const Produs& produs) {

        if(card.buy(produs)) {

            produse.push_back(produs);

            cout << nume
                 << " a cumparat "
                 << produs.getNume()
                 << ".\n";
        }
        else {

            cout << nume
                 << " nu are suficienti bani.\n";
        }
    }

    void afiseazaProduse() {

        cout << "Produsele lui "
             << nume
             << ":\n";

        for(auto produs : produse) {
            cout << produs.getNume() << '\n';
        }
    }

    void afiseazaSold() {
        card.afiseazaSold();
    }
};
\end{lstlisting}

\subsection{Factory Pattern}

O clasă poate fi responsabilă de crearea obiectelor.

Această idee stă la baza design pattern-ului Factory.

\begin{lstlisting}
class Magazin {

public:

    static Produs creeazaLaptop() {
        return Produs("Laptop", 3500);
    }

    static Produs creeazaMouse() {
        return Produs("Mouse", 120);
    }

    static Produs creeazaCasti() {
        return Produs("Casti", 250);
    }
};
\end{lstlisting}

\subsection{Utilizare}

\begin{lstlisting}
int main() {

    Card card("Ion", 5000);

    Persoana persoana("Ion", card);

    Produs laptop = Magazin::creeazaLaptop();
    Produs mouse = Magazin::creeazaMouse();
    Produs casti = Magazin::creeazaCasti();

    persoana.cumpara(mouse);
    persoana.cumpara(casti);
    persoana.cumpara(laptop);

    persoana.afiseazaProduse();
    persoana.afiseazaSold();
}
\end{lstlisting}

\section{Concluzii}

În exemplul anterior am ascuns complet variabila:

\begin{lstlisting}
int bani;
\end{lstlisting}

și am expus doar comportamentul relevant:

\begin{lstlisting}
card.buy(produs)
\end{lstlisting}

Această abordare este mai apropiată de spiritul programării orientate pe obiecte: obiectele expun comportamente și ascund detaliile interne ale implementării.

\end{document}